\chapter{Originality and Future Work}
\begin{spacing}{1.3}

\section{Originality}

Short dummy text goes here. Explain how this project differs from prior work or standard baselines.

\section{Future Work}

One possible direction for future work would be to replace the scikit-learn \texttt{MLPClassifier} with a deep learning
framework such as PyTorch or Keras. While scikit-learn’s MLP is simple and sufficient for a baseline, it lacks several
important training features. It does not support dropout, a regularization technique that randomly deactivates neurons 
during training to reduce overfitting and improve generalization. It also lacks batch normalization, which stabilizes 
training by keeping activations consistent across layers — something that becomes increasingly important in deeper architectures.
Additionally, scikit-learn does not allow per-class or per-sample loss weighting, meaning the model cannot be directed to penalize
false negatives more heavily than false positives. In a cybersecurity context, this is a significant limitation since missing an
attack is typically far more costly than a false alarm. Switching to PyTorch or Keras would address all of these issues and open
the door to custom loss functions such as focal loss, which focuses training on hard-to-classify examples, as well as more advanced
learning rate strategies such as cosine annealing or warm restarts.

This is particularly relevant given that Fuzzers were the hardest attack type to detect across all models in this study, achieving
only 44.75\% recall in MLP and around 62\% in both tree-based models. A more expressive and configurable neural network trained with
focal loss could force the model to focus on these ambiguous samples rather than the easier categories it already handles well. 
Oversampling the Fuzzers class could help improve performance further. Interestingly, Logistic Regression achieved the highest
Fuzzers recall (95.91), suggesting the class may be more linearly separable than expected and that more complex models may be overfitting.
Future work could explore using Logistic Regression specifically for Fuzzers detection as part of an ensemble model.

Another promising direction is model ensembling or stacking. In this study, Logistic Regression, Random Forest, XGBoost, and MLP were
each evaluated independently, yet each model showed different strengths across attack categories. For example, Logistic Regression 
achieved the highest Fuzzers recall while tree-based models excelled on other categories. Rather than selecting a single best model,
a stacking approach could combine the predictions of all four models using a meta-learner, allowing the ensemble to exploit the 
complementary strengths of each classifier.

Finally, all models in this study were trained and evaluated exclusively on the UNSW-NB15 dataset, which was collected in a controlled
lab environment in 2015. This raises questions about how well these models would generalize to real-world traffic or more recent datasets
such as CIC-IDS2017. Since network traffic and attack patterns evolve over time, cross-dataset validation would be a valuable next step 
to better assess the true generalizability of these models.

\end{spacing}
\newpage
